\documentclass[a4paper]{article}

\usepackage[T1]{fontenc}
\usepackage[utf8]{inputenc}
\usepackage[english]{babel}


\usepackage[ruled,vlined,linesnumbered]{algorithm2e} %,noend


% ===== Graph =====
\usepackage{tikz}
% ===== Multicol =====
\usepackage{blindtext}
\usepackage{multicol}
% ===== Cancel =====
\usepackage{cancel}
% ===== Code =====
\usepackage{listings}
\lstdefinestyle{mystyle}{
    breakatwhitespace=false,
    captionpos=b,
    keepspaces=true,
    numbers=left,
    numbersep=5pt,
    showspaces=false,
    showstringspaces=false,
    showtabs=False,
    tabsize=2
}

\usepackage{hyperref}
% \hypersetup{
%   colorlinks=true,
%   linkcolor=black,
%   filecolor=black,
%   urlcolor=black,
% }

\textwidth=450pt
\oddsidemargin=0pt
\textheight=665pt
\voffset=-50pt

\lstset{style=mystyle}

\usepackage{setspace}

\usepackage{amssymb}
\usepackage{amsthm}
\usepackage{mathtools}
\usepackage{bm}

\singlespacing
% ===================
\mathtoolsset{showonlyrefs}

\AtBeginDocument{\renewcommand\proofname{Proof}}

\newcommand{\pluseq}{\mathrel{{+}{=}}}

\newtheorem{theorem}{Theorem}
\newtheorem{corollary}{Corollary}
\newtheorem{lemma}{Lemma}
\newtheorem{remark}{Remark}
\newtheorem{definition}{Definition}

\setcounter{secnumdepth}{3}
\setcounter{tocdepth}{3}

% \title{Statistical Methods for Maching Learning}
% \author{Federico Bruzzone \\ Andrea Longoni \\ Luca Bellani}
% \date{}

\makeindex
\begin{document}

% \maketitle


\begin{titlepage}
    \begin{center}
        % \includegraphics[height=2.6cm]{minerva.pdf} %{Unilogo.png}
        \includegraphics[height=5cm]{minerva.pdf} %{Unilogo.png}

        \vspace*{1.75cm}

        \LARGE
        % \text{University of Milan}

        % \vspace*{1cm}

        \textbf{MSc in Computer Science} \\
        at University of Milan

        \vspace*{1cm}


        \huge
        Formal Methods\\

        \large course held by \textbf{Momigliano Alberto Davide Adolfo}


        \normalsize
        \vspace*{4cm}

        \begin{minipage}[t]{0.47\textwidth}
	       {Email: } \vspace{0.3em} \\
              {\large \href{federico.bruzzone@studenti.unimi.it}{federico.bruzzone@studenti.unimi.it}} \vspace{1em}  \\
        \end{minipage}
        \hfill
        \begin{minipage}[t]{0.47\textwidth}\raggedleft
	       {Created by:} \hspace{-0.9em} \vspace{0.3em} \\
              {\large \textbf{Federico Bruzzone}} \\
              \vspace{1em}
              % {\footnotesize mat. 123456}
        \end{minipage}

        \vfill
        Academic year of 2023/2024

    \end{center}
\end{titlepage}

% \clearpage\null\newpage

% \newpage
% \setlength{\parskip}{0.15em}
% \tableofcontents
\setlength{\parindent}{0pt}
\setlength{\parskip}{0.8em}


% \clearpage\null\newpage

\section{Commutative sum in \(\mathbb{N}\)}

\begin{lemma}\label{lemma:commutative_sum_zero}
  \(\forall \ n \in \mathbb{N}\) it holds that \(n + 0 = n\).
\end{lemma}

\begin{proof}
  We can proof this by induction on \(n\).

  The base case is \(n = 0\), and we have to show that \(0 + 0 = 0\). This is trivially true.

  The inductive step is to show that \(\forall \ n \in \mathbb{N}\) it holds that \(n + 0 = n\).
  We can show that \((n + 1) + 0 = n + 1\) by using the inductive hypothesis.
  \begin{align*}
    (n + 1) + 0 &= n + (1 + 0) \quad \text{by associativity} \\
                &= n + 1 \quad \text{by the inductive hypothesis}
  \end{align*}
\end{proof}

\begin{lemma}\label{lemma:associative_sum}
  \(\forall \ n, m \in \mathbb{N}\) it holds that \(n + (m + 1) = (n + m) + 1\).
\end{lemma}

\begin{proof}
  We can proof this by induction on \(n\).

  The base case is \(n = 0\), and we have to show that \(\forall \ m \in \mathbb{N}\) it holds that \(0 + (m + 1) = (0 + m) + 1\). This is trivially true.

  The inductive step is to show that \(\forall \ n \in \mathbb{N}\) it holds that \(n + (m + 1) = (n + m) + 1\).
  We can show that \((n + 1) + (m + 1) = ((n + 1) + m) + 1\) by using the inductive hypothesis.
  \begin{align*}
    (n + 1) + (m + 1) &= (n + (m + 1)) + 1 \quad \text{by associativity} \\
                      &= ((n + m) + 1) + 1 \quad \text{by the inductive hypothesis} \\
                      &= ((n + 1) + m) + 1 \quad \text{by associativity}
  \end{align*}

\end{proof}

\begin{theorem}\label{theorem:commutative_sum}
\(\forall \ n, m \in \mathbb{N}\) it holds that \(m + n = n + m\).
\end{theorem}

\begin{proof}
  We can proof this by induction on \(n\).

  The base case is \(n = 0\), and we have to show that \(\forall \ m \in \mathbb{N}\) it holds that \(m + 0 = 0 + m\).

  \begin{align*}
    m + 0 &= m \quad \text{by Lemma \ref{lemma:commutative_sum_zero}} \\
  \end{align*}

  The inductive step is to show that \(\forall \ n \in \mathbb{N}\) it holds that \(m + n = n + m\). We can show that \(m + (n + 1) = (n + 1) + m\) by using the inductive hypothesis.

  \begin{align*}
    (n + 1) + m &= (n + m) + 1
  \end{align*}

  \begin{align*}
    m + (n + 1) &= (m + n) + 1 \quad \text{by Lemma \ref{lemma:associative_sum}} \\
                &= (n + m) + 1 \quad \text{by the inductive hypothesis} \\
  \end{align*}

\end{proof}

% \newpage
% \input{chapters/1_introduction.tex}
% \input{chapters/2_nearest_neighbour.tex}
% \input{chapters/3_tree_predictors.tex}
% \input{chapters/4_statistical_learning.tex}
% \input{chapters/5_risk_analisys_for_tree_predictors.tex}
% \input{chapters/6_hyperparameter_tuning_and_risk_estimates.tex}
% \input{chapters/7_consistency_and_nonparametric_algorithms.tex}
% \input{chapters/8_risk_analysis_for_nearest_neighbot.tex}
% \input{chapters/9_linear_predictors.tex}
% \input{chapters/10_online_gradient_descent.tex}
% \input{chapters/11_kernel_functions.tex}
% \input{chapters/12_support_vector_machine.tex}
% \input{chapters/13_stability_and_risk_control_for_SVM.tex}
% \input{chapters/14_neural_networks_and_deep_learning.tex}
% \input{chapters/15_logistic_regression_surrogate_losses.tex}
% \input{chapters/copilot-questions.tex}

% \section{Test Graph}

% \begin{tikzpicture}[node distance={15mm}, thick, main/.style = {draw, circle}]
% \node[main] (1) {$x_1$};

% \node[main] (2) [above right of=1] {$x_2$};
% \node[main] (3) [below right of=1] {$x_3$};
% \node[main] (4) [above right of=3] {$x_4$};
% \node[main] (5) [above right of=4] {$x_5$};
% \node[main] (6) [below right of=4] {$x_6$};

% \draw (2) -- (4);
% \draw[->] (1) -- (2);
% \draw[->] (1) -- (3);

% \draw (1) to (5);
% \draw (1) to [out=135,in=90,looseness=1.5] (5);

% \draw (1) to [out=180,in=270,looseness=5] (1);

% \draw[->] (6) -- node[midway, above right, sloped, pos=1] {+1} (4);

% \end{tikzpicture}

% \begin{tikzpicture}[node distance={15mm}, thick, main/.style = {draw, circle}]
% \node[main] (1) {$x_1$};
% \node[main] (2) [above right of=1] {$x_2$};
% \node[main] (3) [below right of=1] {$x_3$};
% \node[main] (4) [above right of=3] {$x_4$};
% \node[main] (5) [above right of=4] {$x_5$};
% \node[main] (6) [below right of=4] {$x_6$};
% \draw[->] (1) -- (2);
% \draw[->] (1) -- (3);
% \draw (1) to [out=135,in=90,looseness=1.5] (5);
% \draw (1) to [out=180,in=270,looseness=5] (1);
% \draw (2) -- (4);
% \draw (3) -- (4);
% \draw (5) -- (4);
% \draw[->] (5) to [out=315, in=315, looseness=2.5] (3);
% \draw[->] (6) -- node[midway, above right, sloped, pos=1] {+1} (4);
% \end{tikzpicture}

\end{document}
